\section{Собственные значения и собственные векторы}

\begin{definition}
Ненулевой вектор $x$ — \emph{собственный} для оператора $\mathcal A$, если
$\mathcal Ax=\lambda x$ при некотором $\lambda$; число $\lambda$ — \emph{собственное значение}.
Геометрически $\mathcal A$ лишь растягивает (в $\lambda$ раз) направление $x$, не поворачивая его.
\end{definition}

\begin{definition}
\emph{Характеристический многочлен} $\chi_{\mathcal A}(\lambda)=\det(A-\lambda I)$. Собственные
значения — его корни: $\det(A-\lambda I)=0$ (тогда $A-\lambda I$ вырождена и уравнение
$(A-\lambda I)x=0$ имеет ненулевое решение).
\end{definition}

\begin{theorem}[инвариантность $\chi$]
Характеристический многочлен не зависит от выбора базиса. Его коэффициенты при $\lambda^{n-1}$ и
свободный член дают $\tr A$ и $\det A$: $\chi(\lambda)=(-1)^n\lambda^n+(-1)^{n-1}(\tr A)\lambda^{n-1}+\dots+\det A$.
\end{theorem}

\begin{proof}
Для подобной матрицы $A'=S^{-1}AS$:
\[
  \det(A'-\lambda I)=\det\!\bigl(S^{-1}(A-\lambda I)S\bigr)=\det S^{-1}\det(A-\lambda I)\det S=\det(A-\lambda I).
\]
Значит $\chi$ — инвариант. Подставив $\lambda=0$: $\chi(0)=\det A$; коэффициент при
$\lambda^{n-1}$ равен $(-1)^{n-1}\tr A$ (сумма диагональных миноров $1$-го порядка).
\end{proof}

\section{Собственные подпространства}

\begin{theorem}
Множество $V_\lambda=\Ker(\mathcal A-\lambda\mathcal I)=\{x:\mathcal Ax=\lambda x\}$ — подпространство
(\emph{собственное подпространство}). Его размерность называют \emph{геометрической кратностью}
$k_{\text{геом}}(\lambda)$.
\end{theorem}

\begin{proof}
$V_\lambda$ — ядро линейного оператора $\mathcal A-\lambda\mathcal I$, а ядро всегда является
подпространством.
\end{proof}

\begin{theorem}
Собственные векторы, отвечающие \emph{попарно различным} собственным значениям, линейно независимы.
\end{theorem}

\begin{proof}
Индукция по числу $m$ векторов. База $m=1$: один ненулевой вектор независим. Шаг: пусть
$x_1,\dots,x_m$ отвечают различным $\lambda_1,\dots,\lambda_m$ и
$\sum_{i=1}^m c_i x_i=0$. Применим $\mathcal A$ и вычтем уравнение, умноженное на $\lambda_m$:
\[
  0=\mathcal A\Bigl(\sum c_i x_i\Bigr)-\lambda_m\sum c_i x_i=\sum_{i=1}^{m-1}c_i(\lambda_i-\lambda_m)x_i .
\]
По предположению индукции $x_1,\dots,x_{m-1}$ независимы, а $\lambda_i-\lambda_m\ne0$, поэтому все
$c_i=0$ ($i<m$). Тогда $c_m x_m=0$, откуда $c_m=0$.
\end{proof}

\begin{corollary}
Если $\chi$ имеет $n=\dim V$ различных корней, то существует базис из собственных векторов.
\end{corollary}

\section{Алгебраическая и геометрическая кратности}

\begin{definition}
\emph{Алгебраическая кратность} $k_{\text{алг}}(\lambda)$ — кратность $\lambda$ как корня $\chi$.
\emph{Геометрическая} $k_{\text{геом}}(\lambda)=\dim V_\lambda$.
\end{definition}

\begin{theorem}
$1\le k_{\text{геом}}(\lambda)\le k_{\text{алг}}(\lambda)$.
\end{theorem}

\begin{proof}
$\lambda$ — корень $\chi$, значит $V_\lambda\ne\{0\}$ и $k_{\text{геом}}\ge1$. Возьмём базис
$e_1,\dots,e_g$ в $V_\lambda$ ($g=k_{\text{геом}}$) и дополним до базиса $V$. В нём матрица имеет
блочный вид $\begin{psmallmatrix}\lambda I_g & *\\ 0 & B\end{psmallmatrix}$, поэтому
$\chi(\mu)=(\lambda-\mu)^g\det(B-\mu I)$ делится на $(\lambda-\mu)^g$. Значит кратность корня
$\lambda$ не меньше $g$: $k_{\text{алг}}\ge g=k_{\text{геом}}$.
\end{proof}

\begin{example}
Жорданова клетка $\begin{psmallmatrix}2&1\\0&2\end{psmallmatrix}$: $k_{\text{алг}}(2)=2$, но
$V_2=\Span(e_1)$, $k_{\text{геом}}(2)=1$ — кратности не совпадают (матрица не диагонализуема).
\end{example}

\section{Диагонализуемость}

\begin{definition}
Оператор \emph{диагонализуем}, если существует базис из собственных векторов (в нём матрица
диагональна).
\end{definition}

\begin{theorem}[критерий]
Оператор диагонализуем $\iff$ для каждого собственного значения
$k_{\text{геом}}=k_{\text{алг}}$ и сумма алгебраических кратностей равна $n$ (т.е.\ $\chi$
разлагается на линейные множители, и геометрические кратности «добирают» размерность). Равносильно:
$V=\bigoplus_\lambda V_\lambda$.
\end{theorem}

\begin{proof}[Идея]
Собственные подпространства, отвечающие разным $\lambda$, дают независимые векторы (предыдущая
теорема), поэтому $\sum\dim V_\lambda\le n$ с равенством $\iff$ их объединённые базисы образуют
базис всего $V$. Это и есть существование собственного базиса; равенство достигается ровно когда
$\sum k_{\text{геом}}=\sum k_{\text{алг}}=n$.
\end{proof}

\begin{remark}[алгоритм]
1) Найти корни $\chi(\lambda)=0$. 2) Для каждого $\lambda$ решить $(A-\lambda I)x=0$ — базис
$V_\lambda$. 3) Если суммарно набралось $n$ независимых векторов, составить из них столбцы матрицы
$S$; тогда $S^{-1}AS=\operatorname{diag}(\lambda_1,\dots,\lambda_n)$.
\end{remark}
